\subsection{Тестирование}

Для тестирования \glsdisp{sw}{программного обеспечения} (ПО) часто применяются следующие подходы:

\begin{itemize}
    \item \gls{unit-testing} -- проверяется корректность работы отдельных функций или методов изолированно от других компонентов \cite{googletest_documentation,catch2_documentation,boost_test_documentation}.
    \item \gls{module-testing} -- проверяется работа отдельных модулей. Тестируются корректность взаимодействий внутри модуля и его интерфейсов.
    \item \gls{end-to-end-testing} -- решение тестируется полностью на целевых сценариях и данных.
\end{itemize}

Для юнит- и модульного тестирования решения используются фреймворк GoogleTest \cite{googletest_documentation} и тестовая инфраструктура используемой системы сборки CMake \cite{cmake_ctest_documentation,hoffman_martin_cmake}. Также возможно добавление сквозных тестов с использованием реальных или синтетических трасс и инфраструктуры проекта.

Дизайн ключевых компонентов решения позволяет применять юнит- и модульное тестирование, возможно, после нескольких незначительных изменений. Основной задачей для поддержки подобного тестирования является отделение компонентов решения от библиотеки для чтения трасс. Это позволит осуществлять частичную обработку тестовых данных без наличия реальной трассы. Примером подобного отделения является публичный интерфейс модуля ГПД (см. рис. \ref{fig:dfg_uml}): класс \textit{DfgCollector} принимает данные, не привязанные явно к библиотеке для чтения трасс (хотя формат данных сильно похож на формат записей трассы). Это позволяет собирать ГПД и вспомогательные данные на тестовых сценариях, а затем проверять консистентность, инварианты и выполнение тестовых условий. Для упрощения описания тестовых сценариев была реализована вспомогательная инфраструктура, в упрощенном виде представленная в листингах \ref{lst:testing_dfg_builder} и \ref{lst:testing_dfg_collector}. Пример юнит-теста, использующего описанную инфраструктуру, представлен в листинге \ref{lst:testing_dfg_collector_example}.

С помощью описанной тестовой инфраструктуры было добавлено 7 юнит-тестов, покрывающих базовый функционал для работы с ГПД и ряд краевых случаев. Примененный подход показал свою эффективность: тестовые сценарии описываются легко и могут быть повторно использованы несколькими тестами для проверки различных компонентов. Полное тестирование модуля ГПД и других модулей может быть осуществлено в рамках дальнейшего развития решения.

\begin{listing}[htbp]
    \inputminted{c++}{listings/testing_dfg_builder.cpp}
    \caption{Класс для генерации тестовых ГПД.}
    \label{lst:testing_dfg_builder}
\end{listing}

\begin{listing}[htbp]
    \inputminted{c++}{listings/testing_dfg_collector.cpp}
    \caption{Класс-фикстура для юнит-тестирования класса DfgCollector.}
    \label{lst:testing_dfg_collector}
\end{listing}

\begin{listing}[htbp]
    \inputminted{c++}{listings/testing_dfg_collector_example.cpp}
    \caption{Пример юнит-теста для класса DfgCollector.}
    \label{lst:testing_dfg_collector_example}
\end{listing}

\FloatBarrier
